%! AP Statistics — Unit 3, Lesson 3.7 — Lesson Plan
\documentclass[10pt]{article}
\usepackage{apstats-article}
\usepackage{apstats-boxes}

\newcommand{\UnitNumberName}{Unit 3: Collecting Data \quad}
\newcommand{\LessonNumberName}{Lesson 3.7: Inference and Experiments}

\begin{document}
\begin{center}
    {\Large\bfseries \CourseName: \SchoolYear\ (\MeetingLength) \\
    {\small\bfseries \UnitNumberName \LessonNumberName}}
\end{center}

\begin{tcolorbox}[colback=sky, colframe=navy, boxrule=0.9pt, arc=2mm,
    left=2mm, right=2mm, top=1.5mm, bottom=1.5mm]
    \textbf{Primary Objective:} Students will interpret the results of a well-designed
    experiment, determine whether observed differences are statistically significant,
    and apply the scope of inference (causation and generalizability) based on how
    the study was designed (Big Idea: VAR).
\end{tcolorbox}

\renewcommand{\arraystretch}{1.6}
\begin{skillbox}[Priority Ideas \& Skills]{goldbox}
\begin{minipage}[t]{0.48\linewidth}
    \textbf{AP Skill 4 --- Statistical Argumentation}
    \begin{itemize}
        \item Interpret the results of a well-designed experiment (VAR-3.E, Skill 4.B).
        \item Apply statistical inference: random assignment supports causal conclusions
              (VAR-3.E.2, VAR-3.E.3).
        \item Determine generalizability based on whether random selection was used
              (VAR-3.E.4).
    \end{itemize}
\end{minipage}
\hfill
\begin{minipage}[t]{0.48\linewidth}
    \textbf{Key Understandings}
    \begin{itemize}
        \item Statistical significance: observed differences too large to be explained
              by chance alone.
        \item Random \emph{assignment} $\Rightarrow$ can conclude \emph{causation}.
        \item Random \emph{selection} $\Rightarrow$ can \emph{generalize} to the
              population.
        \item Both random assignment AND random selection $\Rightarrow$ can conclude
              causation and generalize.
        \item Neither $\Rightarrow$ can only describe the sample.
    \end{itemize}
\end{minipage}
\end{skillbox}

\begin{skillbox}[{Vocabulary, Concepts, \& Theorems}]{greenbox}
\begin{minipage}[t]{0.48\linewidth}
    \begin{itemize}
        \item \textbf{Statistical inference} --- using data from a study to draw
              conclusions about the population or treatment
        \item \textbf{Statistically significant} --- a result so large it is unlikely
              to have occurred by chance alone
        \item \textbf{Scope of inference} --- what conclusions (causation,
              generalization) are valid given the study design
    \end{itemize}
\end{minipage}
\hfill
\begin{minipage}[t]{0.48\linewidth}
    \begin{itemize}
        \item \textbf{Causal conclusion} --- concluding that the treatment \emph{caused}
              the observed effect; requires random assignment
        \item \textbf{Generalization} --- applying results to the broader population;
              requires random selection
        \item \textbf{Scope of inference table} --- a $2 \times 2$ table summarizing
              what can be concluded based on random assignment and random selection
    \end{itemize}
\end{minipage}
\end{skillbox}

\begin{skillbox}[Activate Prior Knowledge \& Spiral Review]{sky}
\begin{minipage}[t]{0.48\linewidth}
    \textbf{Prior Knowledge Check (3--4 min)}
    \begin{itemize}
        \item Recall: ``What is the difference between random selection and random
              assignment?''
        \item Recall: ``Can an observational study establish causation?'' (No ---
              Lesson 3.2)
        \item Transition: ``Today we build the complete decision framework: what can
              we conclude, and about whom?''
    \end{itemize}
\end{minipage}
\hfill
\begin{minipage}[t]{0.48\linewidth}
    \textbf{Connections Forward}
    \begin{itemize}
        \item Formal significance testing (p-values) $\to$ Units 6--9
        \item Using randomization distributions $\to$ Unit 4 (simulation)
        \item Causal language in AP free-response $\to$ throughout the course
    \end{itemize}
\end{minipage}
\end{skillbox}

\begin{skillbox}[Hook \& Lesson]{sky}
\begin{minipage}[t]{0.48\linewidth}
    \textbf{Hook (5--7 min)}\\
    Display four study results:
    \begin{enumerate}
        \item Survey of 50 volunteers; found a positive association.
        \item Random sample of 1{,}000; found a positive association.
        \item Experiment with random assignment of 80 subjects; found a significant
              difference.
        \item Experiment with random assignment from a random sample; found a
              significant difference.
    \end{enumerate}
    Ask: ``For each, can we conclude causation? Can we generalize to all adults?''
\end{minipage}
\hfill
\begin{minipage}[t]{0.48\linewidth}
    \textbf{Lesson Flow (15--18 min)}
    \begin{enumerate}
        \item Statistical inference: conclusions attributed to the distribution from
              which data were collected (VAR-3.E.1).
        \item Random assignment allows causal conclusions: the treatments, not
              confounding variables, caused the difference (VAR-3.E.2, .E.3).
        \item Statistical significance: the observed difference is too large to be
              due to chance alone.
        \item Generalizability: requires random selection from the population of
              interest (VAR-3.E.4).
        \item Scope of inference table: $2\times2$ grid with random assignment
              (yes/no) vs.\ random selection (yes/no).
    \end{enumerate}
\end{minipage}
\end{skillbox}

\begin{skillbox}[Explicit Instruction \& Active Monitoring]{sky}
\begin{minipage}[t]{0.48\linewidth}
    \textbf{Direct Instruction (10 min)}
    \begin{itemize}
        \item Build the scope-of-inference table cell by cell with student input.
        \item Worked example: a researcher randomly selects 200 adults and randomly
              assigns them to a new exercise program or a control. After 6 months,
              the exercise group loses significantly more weight. Conclusion: the
              program causes weight loss AND we can generalize to all adults.
        \item Counter-example: volunteers randomly assigned to treatments. Can conclude
              causation but NOT generalize beyond similar volunteers.
    \end{itemize}
\end{minipage}
\hfill
\begin{minipage}[t]{0.48\linewidth}
    \textbf{Active Monitoring}
    \begin{itemize}
        \item Listen for: students saying ``the experiment proves the drug works for
              everyone'' when the sample was not randomly selected.
        \item Cold-call: ``If we randomly assign but don't randomly select, what's the
              limitation of our conclusion?'' (Can say treatment caused the effect in
              these subjects, but can't generalize.)
        \item Check scope-of-inference table entries for accuracy.
    \end{itemize}
\end{minipage}
\end{skillbox}

\begin{skillbox}[Group Work \& Differentiation]{redbox}
\begin{minipage}[t]{0.48\linewidth}
    \textbf{Partner Activity (8--10 min)}\\
    Four study summaries (observational and experimental, random and non-random
    selection). Partners:
    \begin{enumerate}
        \item Identify whether random assignment was used.
        \item Identify whether random selection was used.
        \item State the appropriate causal conclusion (or lack thereof).
        \item State to whom the results can be generalized.
    \end{enumerate}
\end{minipage}
\hfill
\begin{minipage}[t]{0.48\linewidth}
    \textbf{Differentiation}
    \begin{itemize}
        \item \textit{Support}: Scope of inference reference card with the $2 \times 2$
              table filled in.
        \item \textit{Extension}: Write four different conclusions for the same data
              set, one for each cell of the scope table, and explain what would need
              to change in the study design to support each.
        \item \textit{ELL}: Color-coded table (green = yes, red = no for each type of
              randomization).
    \end{itemize}
\end{minipage}
\end{skillbox}

\begin{skillbox}[Individual Work \& Assessment]{redbox}
\begin{minipage}[t]{0.48\linewidth}
    \textbf{Exit Ticket (5 min)}\\
    Scenario: \textit{``A school counselor randomly selects 60 students from all
    students in the district and randomly assigns them to either a mindfulness
    program or study hall. After 8 weeks, the mindfulness group has significantly
    lower stress scores.''}
    \begin{enumerate}
        \item Was random assignment used?
        \item Was random selection used?
        \item Can we conclude the program caused lower stress?
        \item Can we generalize to all students in the district?
    \end{enumerate}
\end{minipage}
\hfill
\begin{minipage}[t]{0.48\linewidth}
    \textbf{AP-Style Check}\\
    \textit{``Researchers randomly assigned 120 patients to either a new medication
    or a placebo. The medication group showed significantly greater improvement.
    Patients were recruited from one hospital. Which conclusion is best
    supported?''}
    \begin{enumerate}[label=(\Alph*)]
        \item The medication causes improvement and this generalizes to all patients.
        \item The medication caused improvement in these patients, but results may
              not generalize to all patients.
        \item There is an association, but causation cannot be established.
        \item Results can be generalized but causation cannot be established.
    \end{enumerate}
    Answer: (B)
\end{minipage}
\end{skillbox}

\begin{skillbox}[Reinforcement \& Extension]{goldbox}
\begin{minipage}[t]{0.48\linewidth}
    \textbf{Homework / Practice}
    \begin{itemize}
        \item For 4 study descriptions: identify randomization type(s), state the
              appropriate scope of inference, and write a proper conclusion sentence.
        \item Reflection: \textit{``A study uses random assignment but not random
              selection. What can and cannot be concluded?''}
    \end{itemize}
\end{minipage}
\hfill
\begin{minipage}[t]{0.48\linewidth}
    \textbf{Extension / Preview}
    \begin{itemize}
        \item \textit{Extension}: Find three AP free-response questions from past exams
              involving experiments. For each, identify whether the conclusion drawn in
              the scoring rubric reflects causation, generalization, or both.
        \item \textit{Preview}: Unit 4 introduces probability --- the mathematical
              foundation for ``statistically significant'' used throughout Unit 3.
    \end{itemize}
\end{minipage}
\end{skillbox}

\end{document}
